\documentclass{natsirt}

\usepackage{pgfplots}
\pgfplotsset{compat=1.18}

\title{Calculus: Concept Checklist}
\author{Calculus I (Y1S1)}
\subtitle{Active-Recall \& Self-Testing Revision Guide}

\setlength{\parindent}{0pt}
\setlength{\parskip}{3pt}
\setlength{\headheight}{15pt}

% Custom Box Environments for Pure Active-Recall Checklist
\newcounter{topic}[section]

\newenvironment{recallbox}[1][]{%
    \begin{bluebox}[#1]\refstepcounter{topic}\textbf{Definitions \& Concepts to State \thesection.\thetopic}%
}{\end{bluebox}}

\newenvironment{theoremsbox}[1][]{%
    \begin{greenbox}[#1]\refstepcounter{topic}\textbf{Theorems \& Identities to State \thesection.\thetopic}%
}{\end{greenbox}}

\newenvironment{challengebox}[1][]{%
    \begin{redbox}[#1]\refstepcounter{topic}\textbf{Proof Challenges \& Calculations \thesection.\thetopic}%
}{\end{redbox}}

\newenvironment{warningbox}[1][]{\begin{whitebox}[#1]\textbf{Conventions \& Exam Pitfalls. }}{\end{whitebox}}
\newenvironment{reviewbox}[1][]{\begin{whitebox}[#1]\textbf{Textbook Exercise Checklist. }}{\end{whitebox}}

\newcommand{\citem}{\item[$\square$]}
\newcommand{\citemdone}{\item[$\text{\rlap{$\checkmark$}}\square$]}
\newcommand{\Img}{\operatorname{Im}}

\begin{document}

% =========================================================================
% CHAPTER 1: FUNCTIONS, MODELS, AND TRANSFORMATIONS
% =========================================================================
\section{Functions, Models, and Transformations}

\begin{warningbox}
    \begin{itemize}
        \citem \textbf{Horizontal vs. Vertical Scaling:} For $c > 1$, $f(cx)$ is a horizontal \emph{compression} by factor $c$, whereas $c f(x)$ is a vertical \emph{stretch} by factor $c$.
        \citem \textbf{Trigonometric Periods:} The fundamental period of $\sin x$ is $2\pi$, while the period of $\sin(ax)$ ($a > 0$) is compressed to $\frac{2\pi}{a}$.
    \end{itemize}
\end{warningbox}

\begin{recallbox}
    \begin{itemize}
        \citem \textbf{Domain and Range:} Define the domain $D$ and image (range) $\Img(f) = \set{f(x) : x \in D}$ of a real-valued function $f$.
        \citem \textbf{Monotonicity:} Define strictly increasing, strictly decreasing, and strictly monotonic functions.
        \citem \textbf{Function Classifications:} Define the following standard families of functions:
        \begin{itemize}
            \item Power functions $x^\alpha$ (polynomial for $\alpha = n \in \Z^+$, root for $\alpha = 1/n$, reciprocal for $\alpha = -1$).
            \item Rational functions $R(x) = P(x)/Q(x)$ where $P, Q$ are polynomials.
            \item Algebraic functions (functions constructed using finite algebraic operations and roots).
            \item Periodic functions ($f(x + p) = f(x)$ for all $x$, fundamental period $p > 0$).
            \item Exponential functions $f(x) = a^x$ ($a > 0, a \neq 1$).
        \end{itemize}
        \citem \textbf{Horizontal Line Test:} State the Horizontal Line Test for function injectivity (one-to-one property).
        \citem \textbf{Graph Transformations:} For $c > 0$, state the geometric effect on $y = f(x)$ of:
        \begin{align*}
            y = c f(x), \quad y = \frac{1}{c} f(x), \quad y = f(cx), \quad y = f\pr{\frac{x}{c}}.
        \end{align*}
    \end{itemize}
\end{recallbox}

\begin{theoremsbox}
    \begin{itemize}
        \citem \textbf{Composite Periodicity:} If $f$ has fundamental period $p$, then $g(x) = f(ax)$ has fundamental period $p/a$ for any $a > 0$.
    \end{itemize}
\end{theoremsbox}

\begin{challengebox}
    \begin{itemize}
        \citem \textbf{Period Derivation:} Prove that the fundamental period of $f(x) = \sin(ax)$ is $\frac{2\pi}{a}$ using the definition $f(x + p) = f(x)$.
    \end{itemize}
\end{challengebox}


% =========================================================================
% CHAPTER 2: LIMITS AND CONTINUITY
% =========================================================================
\section{Limits, Epsilon-Delta Proofs, and Continuity}

\begin{warningbox}
    \begin{itemize}
        \citem \textbf{Horizontal Asymptotes Requirement:} When finding horizontal asymptotes, you must evaluate \emph{both} directional limits: $x \to +\infty$ and $x \to -\infty$.
        \citem \textbf{$\epsilon$-$\delta$ Bounding Strategy:} In quadratic limit proofs $\lim_{x\to a} f(x) = L$, choose an initial bound (e.g., $\delta_1 = 1$) to bound the extra factor $|x + c| \le M$, then set $\delta = \min\set{1, \frac{\epsilon}{M}}$.
    \end{itemize}
\end{warningbox}

\begin{recallbox}
    \begin{itemize}
        \citem \textbf{Precise Limit Definition ($\epsilon$-$\delta$):} State the formal $(\eps, \delta)$-definition of $\DS\lim_{x\to a} f(x) = L$:
        \begin{align*}
            \forall \eps > 0, \, \exists \delta > 0 \text{ such that } 0 < |x - a| < \delta \implies |f(x) - L| < \eps.
        \end{align*}
        \citem \textbf{One-Sided Limits:} State the $(\eps, \delta)$-definitions of the left-hand limit $\lim_{x\to a^-} f(x) = L$ and right-hand limit $\lim_{x\to a^+} f(x) = L$.
        \citem \textbf{Infinite Limits \& Asymptotes:}
        \begin{itemize}
            \item Define infinite limit $\lim_{x\to a} f(x) = \infty$ and $\lim_{x\to a} f(x) = -\infty$ (vertical asymptote $x = a$).
            \item Define limits at infinity $\lim_{x\to\infty} f(x) = L$ and $\lim_{x\to-\infty} f(x) = L$ (horizontal asymptote $y = L$).
        \end{itemize}
        \citem \textbf{Continuity Definitions:} State what it means for $f$ to be:
        \begin{itemize}
            \item Continuous at a point $x = a$ ($\lim_{x\to a} f(x) = f(a)$).
            \item Continuous from the left / right at $x = a$.
            \item Continuous on an open interval $(a, b)$ and on a closed interval $[a, b]$.
        \end{itemize}
    \end{itemize}
\end{recallbox}

\begin{theoremsbox}
    \begin{itemize}
        \citem \textbf{Composite Limit Theorem:} If $\lim_{x\to a} g(x) = L$ and $f$ is continuous at $L$, then $\lim_{x\to a} f(g(x)) = f(L)$.
        \citem \textbf{Differentiability Implies Continuity:} If $f$ is differentiable at $x = a$, then $f$ is continuous at $x = a$. (State why the converse is false with $f(x) = |x|$).
        \citem \textbf{Sign-Preserving / Separation Lemma:} If $\lim_{x\to a} g(x) = B \neq 0$, then there exists a positive number $\delta > 0$ such that:
        \begin{align*}
            0 < |x - a| < \delta \implies |g(x)| > \frac{|B|}{2}.
        \end{align*}
    \end{itemize}
\end{theoremsbox}

\begin{challengebox}
    \begin{itemize}
        \citem \textbf{Interval Continuity Proofs:}
        \begin{enumerate}
            \item Show that $f(x) = 1 - \sqrt{1 - x^2}$ is continuous on $[-1, 1]$ (including one-sided endpoints).
            \item Show that $h(x) = |x^2 - 3x - 6|$ is continuous at every real number.
        \end{enumerate}
        \citem \textbf{Differentiability $\implies$ Continuity Proof:} Prove that if $f'(a) = \lim_{x\to a} \frac{f(x)-f(a)}{x-a}$ exists, then $\lim_{x\to a} f(x) = f(a)$.
        \citem \textbf{Separation Lemma Proof:} Prove the Separation Lemma using the formal $\epsilon$-$\delta$ limit definition with $\eps = \frac{|B|}{2}$.
    \end{itemize}
\end{challengebox}

\begin{reviewbox}
    \begin{center}
    \begin{tabular}{|l|l|c|}
        \hline
        \textbf{Section} & \textbf{Textbook Problems} & \textbf{Status} \\
        \hline
        Section 2.4 & Problems 39, 43 & $\square$ Done \\
        \hline
    \end{tabular}
    \end{center}
\end{reviewbox}


% =========================================================================
% CHAPTER 3: DIFFERENTIATION, TRIGONOMETRIC LIMITS, AND DIFFERENTIALS
% =========================================================================
\section{Differentiation, Special Trigonometric Limits, and Differentials}

\begin{recallbox}
    \begin{itemize}
        \citem \textbf{Derivative Definition:} State the definition of the derivative $f'(x) = \lim_{h\to 0} \frac{f(x+h) - f(x)}{h}$.
        \citem \textbf{Chain Rule Formula:} State the Chain Rule for composite functions: $\od{}{x} f(g(x)) = f'(g(x)) g'(x)$.
        \citem \textbf{Differentials vs. Increments:} Distinguish among:
        \begin{itemize}
            \item Independent variable increment $\Delta x$ and differential $dx$ ($dx = \Delta x$).
            \item Dependent variable true increment $\Delta y = f(x + \Delta x) - f(x)$.
            \item Dependent variable differential $dy = f'(x) dx$ and the linear approximation $\Delta y \approx dy$.
        \end{itemize}
    \end{itemize}
\end{recallbox}

\begin{theoremsbox}
    \begin{itemize}
        \citem \textbf{Fundamental Trigonometric Limits:}
        \begin{align*}
            \lim_{x\to 0} \frac{\sin x}{x} = 1 \quad \text{and} \quad \lim_{x\to 0} \frac{1 - \cos x}{x} = 0.
        \end{align*}
        \citem \textbf{Chain Rule Theorem:} If $g$ is differentiable at $x$ and $f$ is differentiable at $g(x)$, then the composite function $F = f \circ g$ is differentiable at $x$ with $F'(x) = f'(g(x))g'(x)$.
    \end{itemize}
\end{theoremsbox}

\begin{challengebox}
    \begin{itemize}
        \citem \textbf{Trigonometric Limits Evaluation:} Evaluate $\lim_{x\to 0} \frac{1 - \cos x}{x}$ algebraically using half-angle identities and the squeeze theorem.
        \citem \textbf{Chain Rule Proof:} Write the full formal proof of the Chain Rule using the differentiability error formulation $\Delta y = f'(a)\Delta u + \eps \Delta u$.
        \citem \textbf{Differential Error Estimation:} Calculate $dy$ versus $\Delta y$ for $y = x^3 + 2x$ as $x$ changes from $2$ to $2.05$.
    \end{itemize}
\end{challengebox}

\begin{reviewbox}
    \begin{center}
    \begin{tabular}{|l|l|c|}
        \hline
        \textbf{Section} & \textbf{Textbook Problems} & \textbf{Status} \\
        \hline
        Section 3.5 & Problem 75 & $\square$ Done \\
        Section 3.9 & Problems 35, 42 & $\square$ Done \\
        Section 3.10& Problems 33, 35, 40 & $\square$ Done \\
        \hline
    \end{tabular}
    \end{center}
\end{reviewbox}


% =========================================================================
% CHAPTER 4: APPLICATIONS OF DIFFERENTIATION
% =========================================================================
\section{Applications of Derivatives, Mean Value Theorem, and Curve Sketching}

\begin{recallbox}
    \begin{itemize}
        \citem \textbf{Extrema and Critical Points:}
        \begin{itemize}
            \item Define absolute (global) maximum/minimum and local (relative) maximum/minimum.
            \item Define a \emph{critical point} ($f'(c) = 0$) and a \emph{singular point} ($f'(c)$ undefined).
        \end{itemize}
        \citem \textbf{Rolle's Theorem \& Mean Value Theorem:} State the hypotheses and conclusions of:
        \begin{itemize}
            \item Rolle's Theorem
            \item Mean Value Theorem (MVT)
        \end{itemize}
        \citem \textbf{Concavity and Inflection Points:} Define concave up, concave down, and inflection point in terms of $f'(x)$ and $f''(x)$.
        \citem \textbf{8-Step Curve Sketching Guide:} List all 8 standard systematic steps for sketching the graph of $y = f(x)$:
        \begin{enumerate}
            \item Domain
            \item Intercepts ($x$ and $y$ intercepts)
            \item Symmetry (Even, Odd, Periodic)
            \item Asymptotes (Horizontal, Vertical, Slant)
            \item Intervals of Increase and Decrease (First Derivative Test)
            \item Local Extrema
            \item Concavity and Inflection Points (Second Derivative Test)
            \item Sketch the Curve
        \end{enumerate}
        \citem \textbf{Newton-Raphson Method:} State the iteration formula for finding roots:
        \begin{align*}
            x_{n+1} = x_n - \frac{f(x_n)}{f'(x_n)}.
        \end{align*}
    \end{itemize}
\end{recallbox}

\begin{theoremsbox}
    \begin{itemize}
        \citem \textbf{Extreme Value Theorem (EVT):} If $f$ is continuous on a closed interval $[a, b]$, then $f$ attains an absolute maximum and an absolute minimum on $[a, b]$.
        \citem \textbf{Fermat's Theorem (Critical Point Theorem):} If $f$ has a local extremum at $c$ and $f'(c)$ exists, then $f'(c) = 0$.
        \citem \textbf{Constant Function Theorem:} If $F'(x) = 0$ for all $x \in (a, b)$, then $F(x) = C$ is constant on $(a, b)$.
        \citem \textbf{Equal Derivatives Theorem:} If $f'(x) = g'(x)$ for all $x \in (a, b)$, then $f(x) = g(x) + C$ on $(a, b)$.
    \end{itemize}
\end{theoremsbox}

\begin{challengebox}
    \begin{itemize}
        \citem \textbf{Uniqueness of Real Root:} Prove that the cubic equation $x^3 + x - 1 = 0$ has \emph{strictly one real root} using the Intermediate Value Theorem (existence) and Rolle's Theorem / MVT (uniqueness).
        \citem \textbf{Zero Derivative Characterization Proof:} Prove that if $F'(x) = 0$ on $(a, b)$, then $F(x) = C$ using the Mean Value Theorem.
        \citem \textbf{Newton-Raphson Derivation:} Derive the Newton-Raphson iteration formula geometrically from the tangent line approximation at $(x_n, f(x_n))$.
    \end{itemize}
\end{challengebox}

\begin{reviewbox}
    \begin{center}
    \begin{tabular}{|l|l|c|}
        \hline
        \textbf{Section} & \textbf{Textbook Problems} & \textbf{Status} \\
        \hline
        Section 4.2 & Problems 19, 25, 29 & $\square$ Done \\
        Section 4.3 & Problems 75, 83 & $\square$ Done \\
        Section 4.4 & Problem 55 & $\square$ Done \\
        Section 4.5 & Problem 73 & $\square$ Done \\
        Section 4.8 & Problem 17 & $\square$ Done \\
        \hline
    \end{tabular}
    \end{center}
\end{reviewbox}


% =========================================================================
% CHAPTER 5: DEFINITE INTEGRALS AND FUNDAMENTAL THEOREMS
% =========================================================================
\section{Definite Integrals and Fundamental Theorems of Calculus}

\begin{recallbox}
    \begin{itemize}
        \citem \textbf{Riemann Sums Definition:} Define the Riemann sum $\sum_{i=1}^n f(x_i^*) \Delta x$ of $f$ on $[a, b]$ where $\Delta x = \frac{b-a}{n}$. Define Left sum ($x_i^* = x_{i-1}$), Right sum ($x_i^* = x_i$), and Midpoint sum ($x_i^* = \bar{x}_i$).
        \citem \textbf{Definite Integral Definition:} Define $\int_a^b f(x) dx = \lim_{n\to\infty} \sum_{i=1}^n f(x_i^*) \Delta x$.
        \citem \textbf{Net Displacement vs. Distance Traveled:}
        \begin{itemize}
            \item Net displacement $= \int_{t_1}^{t_2} v(t) dt$.
            \item Total distance traveled $= \int_{t_1}^{t_2} |v(t)| dt$.
        \end{itemize}
    \end{itemize}
\end{recallbox}

\begin{theoremsbox}
    \begin{itemize}
        \citem \textbf{Integrability of Continuous Functions:} If $f$ is continuous on $[a, b]$, or bounded with only finitely many jump discontinuities, then $f$ is integrable on $[a, b]$.
        \citem \textbf{Comparison and Boundedness Properties:}
        \begin{itemize}
            \item If $f(x) \le g(x)$ on $[a, b]$, then $\int_a^b f(x) dx \le \int_a^b g(x) dx$.
            \item If $m \le f(x) \le M$ on $[a, b]$, then $m(b - a) \le \int_a^b f(x) dx \le M(b - a)$.
        \end{itemize}
        \citem \textbf{First Fundamental Theorem of Calculus (FTC 1):} Let $f$ be continuous on $[a, b]$. If $g(x) = \int_a^x f(t) dt$ for $x \in [a, b]$, then $g$ is continuous on $[a, b]$, differentiable on $(a, b)$, and:
        \begin{align*}
            g'(x) = \od{}{x} \int_a^x f(t) dt = f(x).
        \end{align*}
        \citem \textbf{Second Fundamental Theorem of Calculus (FTC 2):} If $f$ is continuous on $[a, b]$ and $F$ is any antiderivative of $f$ ($F' = f$), then:
        \begin{align*}
            \int_a^b f(x) dx = F(b) - F(a) = \br{F(x)}_a^b.
        \end{align*}
        \citem \textbf{Symmetric Integral Theorem:} For an integrable function $f$ on $[-a, a]$:
        \begin{align*}
            \int_{-a}^a f(x) dx = \cases{
                2\int_0^a f(x) dx & \text{if } f \text{ is even } (f(-x) = f(x)) \\
                0 & \text{if } f \text{ is odd } (f(-x) = -f(x))
            }
        \end{align*}
    \end{itemize}
\end{theoremsbox}

\begin{challengebox}
    \begin{itemize}
        \citem \textbf{FTC 1 Proof:} Write the complete proof of FTC 1 using the definition of derivative and the Extreme Value Theorem / Boundedness property.
        \citem \textbf{FTC 2 Proof:} Prove FTC 2 using FTC 1 and the Constant Function Theorem.
        \citem \textbf{Integral Estimation Task:} Give rigorous upper and lower bounds for $\int_0^1 e^{-x^2} dx$ using the boundedness property on $[0, 1]$.
    \end{itemize}
\end{challengebox}

\begin{reviewbox}
    \begin{center}
    \begin{tabular}{|l|l|c|}
        \hline
        \textbf{Section} & \textbf{Textbook Problems} & \textbf{Status} \\
        \hline
        Section 5.1 & Problem 5 & $\square$ Done \\
        Section 5.2 & Problems 55, 57 & $\square$ Done \\
        Section 5.3 & Problems 13, 17, 75 & $\square$ Done \\
        Section 5.4 & Problems 45, 59, 61 & $\square$ Done \\
        Section 5.5 & Problems 61, 67, 73, 91, 93 & $\square$ Done \\
        \hline
    \end{tabular}
    \end{center}
\end{reviewbox}


% =========================================================================
% CHAPTER 6: APPLICATIONS OF DEFINITE INTEGRALS
% =========================================================================
\section{Applications of Definite Integrals: Volumes and Average Values}

\begin{recallbox}
    \begin{itemize}
        \citem \textbf{Solids of Revolution Methods:}
        \begin{itemize}
            \item \textbf{Disk Method:} $V = \pi \int_a^b [R(x)]^2 dx$ (cross-sections perpendicular to axis of revolution).
            \item \textbf{Washer Method:} $V = \pi \int_a^b \pr{[R(x)]^2 - [r(x)]^2} dx$ (outer radius $R(x)$, inner radius $r(x)$).
            \item \textbf{Cylindrical Shells Method:} $V = 2\pi \int_a^b (\text{shell radius})(\text{shell height}) dx$ (cylindrical shells parallel to axis of revolution).
        \end{itemize}
        \citem \textbf{Average Value of a Function:} State the definition of the average value $f_{\text{avg}}$ of $f$ on $[a, b]$:
        \begin{align*}
            f_{\text{avg}} = \frac{1}{b - a} \int_a^b f(x) dx.
        \end{align*}
    \end{itemize}
\end{recallbox}

\begin{theoremsbox}
    \begin{itemize}
        \citem \textbf{Mean Value Theorem for Integrals:} If $f$ is continuous on $[a, b]$, then there exists a number $c \in [a, b]$ such that:
        \begin{align*}
            f(c) = f_{\text{avg}} = \frac{1}{b - a} \int_a^b f(x) dx \iff \int_a^b f(x) dx = f(c)(b - a).
        \end{align*}
    \end{itemize}
\end{theoremsbox}

\begin{challengebox}
    \begin{itemize}
        \citem \textbf{Washer vs. Shell Setup:} Set up definite integrals to find the volume of the solid generated by revolving the region bounded by $y = x^2$ and $y = 2x$ about:
        \begin{enumerate}
            \item The $x$-axis (Washer method).
            \item The line $x = 3$ (Cylindrical shells method).
        \end{enumerate}
        \citem \textbf{MVT for Integrals Proof:} Prove the Mean Value Theorem for Integrals using the Extreme Value Theorem and the Boundedness Property of definite integrals.
    \end{itemize}
\end{challengebox}

\begin{reviewbox}
    \begin{center}
    \begin{tabular}{|l|l|c|}
        \hline
        \textbf{Section} & \textbf{Textbook Problems} & \textbf{Status} \\
        \hline
        Section 6.1 & Problems 3, 15 & $\square$ Done \\
        Section 6.3 & Problems 5, 7, 17, 19 & $\square$ Done \\
        Section 6.5 & Problems 3, 5, 7 & $\square$ Done \\
        \hline
    \end{tabular}
    \end{center}
\end{reviewbox}


% =========================================================================
% CHAPTER 7: TECHNIQUES OF INTEGRATION & NUMERICAL METHODS
% =========================================================================
\section{Techniques of Integration, Numerical Rules, and Improper Integrals}

\begin{recallbox}
    \begin{itemize}
        \citem \textbf{Integration by Parts:} State the formula $\int u \, dv = uv - \int v \, du$. Outline the Tabular (DI) method.
        \citem \textbf{Trigonometric Product-to-Sum \& Sum-to-Product Formulas:}
        \begin{align*}
            \sin A \cos B &= \frac{1}{2}\br{\sin(A+B) + \sin(A-B)}, \\
            \cos A \cos B &= \frac{1}{2}\br{\cos(A+B) + \cos(A-B)}, \\
            \sin A + \sin B &= 2\sin\pr{\frac{A+B}{2}}\cos\pr{\frac{A-B}{2}}, \\
            \cos A - \cos B &= -2\sin\pr{\frac{A+B}{2}}\sin\pr{\frac{A-B}{2}}.
        \end{align*}
        \citem \textbf{Numerical Integration Rules:}
        \begin{itemize}
            \item Trapezoidal Rule:
            \begin{align*}
                T_n = \frac{\Delta x}{2}\br{f(x_0) + 2f(x_1) + 2f(x_2) + \dots + 2f(x_{n-1}) + f(x_n)}.
            \end{align*}
            \item Simpson's Rule ($n$ even):
            \begin{align*}
                S_n = \frac{\Delta x}{3}\br{f(x_0) + 4f(x_1) + 2f(x_2) + 4f(x_3) + \dots + 4f(x_{n-1}) + f(x_n)}.
            \end{align*}
        \end{itemize}
        \citem \textbf{Improper Integrals Classifications:}
        \begin{itemize}
            \item Type 1 (Infinite intervals): $\int_a^\infty f(x) dx = \lim_{t\to\infty} \int_a^t f(x) dx$.
            \item Type 2 (Discontinuous integrands): If $f$ is discontinuous at $b$, $\int_a^b f(x) dx = \lim_{t\to b^-} \int_a^t f(x) dx$.
        \end{itemize}
    \end{itemize}
\end{recallbox}

\begin{theoremsbox}
    \begin{itemize}
        \citem \textbf{$p$-Integral Convergence Test:}
        \begin{align*}
            \int_1^\infty \frac{1}{x^p} dx \text{ converges} \iff p > 1 \quad \text{and} \quad \int_0^1 \frac{1}{x^p} dx \text{ converges} \iff p < 1.
        \end{align*}
        \citem \textbf{Direct Comparison Test for Improper Integrals:} Let $f(x) \ge g(x) \ge 0$ on $[a, \infty)$.
        \begin{itemize}
            \item If $\int_a^\infty f(x) dx$ converges, then $\int_a^\infty g(x) dx$ converges.
            \item If $\int_a^\infty g(x) dx$ diverges, then $\int_a^\infty f(x) dx$ diverges.
        \end{itemize}
    \end{itemize}
\end{theoremsbox}

\begin{challengebox}
    \begin{itemize}
        \citem \textbf{Product-to-Sum Derivations:} Derive $\sin A \cos B = \frac{1}{2}[\sin(A+B) + \sin(A-B)]$ using the angle addition formulas for sine.
        \citem \textbf{Simpson's Rule Derivation:} Derive Simpson's rule by integrating the interpolating parabola through three equally spaced points $(-\Delta x, y_0), (0, y_1), (\Delta x, y_2)$.
        \citem \textbf{Improper Integral Evaluation:} Determine whether $\int_0^\infty \frac{1}{1 + x^2} dx$ and $\int_0^1 \frac{1}{\sqrt{x}} dx$ converge, and evaluate their exact values.
    \end{itemize}
\end{challengebox}

\begin{reviewbox}
    \begin{center}
    \begin{tabular}{|l|l|c|}
        \hline
        \textbf{Section} & \textbf{Textbook Problems} & \textbf{Status} \\
        \hline
        Section 7.1 & Problems 9, 15, 17, 33, 41, 61, 63, 65 & $\square$ Done \\
        Section 7.2 & Problem 41 & $\square$ Done \\
        Section 7.3 & Problems 13, 29, 37 & $\square$ Done \\
        Section 7.4 & Problems 15, 29, 31, 47, 51 & $\square$ Done \\
        Section 7.5 & Problem 31 & $\square$ Done \\
        Section 7.7 & Problem 7 & $\square$ Done \\
        Section 7.8 & Problems 5, 13, 23, 31, 49, 51 & $\square$ Done \\
        \hline
    \end{tabular}
    \end{center}
\end{reviewbox}


% =========================================================================
% CHAPTER 8: FURTHER APPLICATIONS (ARC LENGTH, SURFACE AREA, CENTROIDS)
% =========================================================================
\section{Further Applications: Arc Length, Surface Area, and Centroids}

\begin{recallbox}
    \begin{itemize}
        \citem \textbf{Arc Length Formulas:}
        \begin{itemize}
            \item Cartesian curve $y = f(x)$ on $[a, b]$: $L = \int_a^b \sqrt{1 + [f'(x)]^2} dx$.
            \item Parametric curve $x = f(t), y = g(t)$ on $[\alpha, \beta]$: $L = \int_\alpha^\beta \sqrt{[f'(t)]^2 + [g'(t)]^2} dt$.
            \item Arc length function: $s(x) = \int_a^x \sqrt{1 + [f'(t)]^2} dt$.
        \end{itemize}
        \citem \textbf{Surface Area of Revolution:}
        \begin{itemize}
            \item Rotation about $x$-axis: $S = 2\pi \int_a^b y \, ds = 2\pi \int_a^b f(x) \sqrt{1 + [f'(x)]^2} dx$.
            \item Rotation about $y$-axis: $S = 2\pi \int_a^b x \, ds = 2\pi \int_a^b x \sqrt{1 + [f'(x)]^2} dx$.
        \end{itemize}
        \citem \textbf{Center of Mass and Centroids:}
        \begin{itemize}
            \item Discrete system of $n$ masses: $\bar{x} = \frac{\sum m_i x_i}{\sum m_i}$.
            \item Planar lamina bounded by $y = f(x)$ and $y = g(x)$ ($f \ge g$):
            \begin{align*}
                A &= \int_a^b [f(x) - g(x)] dx, \\
                \bar{x} &= \frac{1}{A}\int_a^b x[f(x)-g(x)] dx, \\
                \bar{y} &= \frac{1}{A}\int_a^b \frac{[f(x)]^2 - [g(x)]^2}{2} dx.
            \end{align*}
        \end{itemize}
    \end{itemize}
\end{recallbox}

\begin{theoremsbox}
    \begin{itemize}
        \citem \textbf{Theorem of Pappus for Volume:} Let $R$ be a plane region rotated about an external axis in its plane. The volume of the resulting solid is $V = 2\pi \bar{r} A$, where $A$ is the area of $R$ and $\bar{r}$ is the distance from the axis to the centroid of $R$.
        \citem \textbf{Theorem of Pappus for Surface Area:} The area of the surface generated by revolving a curve of length $L$ about an external axis is $S = 2\pi \bar{r} L$, where $\bar{r}$ is the distance from the axis to the centroid of the curve.
    \end{itemize}
\end{theoremsbox}

\begin{challengebox}
    \begin{itemize}
        \citem \textbf{Centroid Derivation for Planar Lamina:} Derive the integral formulas for $\bar{x}$ and $\bar{y}$ by partitioning the region into vertical approximating rectangles of mass $dm = \rho (f(x) - g(x))dx$.
        \citem \textbf{Torus Volume and Surface Area via Pappus:} Use the Theorems of Pappus to find the volume and surface area of a torus obtained by rotating a circle of radius $r$ centered at $(R, 0)$ ($R > r$) about the $y$-axis.
    \end{itemize}
\end{challengebox}

\begin{reviewbox}
    \begin{center}
    \begin{tabular}{|l|l|c|}
        \hline
        \textbf{Section} & \textbf{Textbook Problems} & \textbf{Status} \\
        \hline
        Section 8.1 & Problems 9, 19 & $\square$ Done \\
        Section 8.2 & Problems 9, 13, 17 & $\square$ Done \\
        Section 8.3 & Problems 23, 27, 31, 45 & $\square$ Done \\
        \hline
    \end{tabular}
    \end{center}
\end{reviewbox}

\end{document}
